\section{API y autenticacion}

El nodo expone una API REST sobre HTTP con exactamente cuatro endpoints
funcionales, mas un \codigo{/stats} de monitoreo. Las rutas se declaran en
\codigo{server/BancoHttp.java} y cada una resuelve con un manejador propio.
El registro y el inicio de sesion crean credenciales y emiten un JWT; las
operaciones sobre cuentas y transferencias exigen ese token en cada peticion.

\subsection{Endpoints REST}

Los cuatro endpoints del proyecto son: \codigo{POST /api/register} para crear
una credencial, \codigo{POST /api/login} que devuelve un JWT, \codigo{GET
/api/accounts/\{id\}} para consultar una cuenta y \codigo{POST
/api/transactions/transfer} para mover saldo. Los dos ultimos exigen el header
\codigo{Authorization: Bearer <token>}.

\begin{tablaglab}{L{5.5cm} C{1.6cm} C{1.5cm} L{4cm}}
\thead{Ruta} & \thead{Metodo} & \thead{JWT} & \thead{Manejador} \\ \midrule
/api/register                  & POST & No  & AuthHandler \\
/api/login                     & POST & No  & AuthHandler \\
/api/accounts/\{id\}           & GET  & Si  & AccountHandler \\
/api/transactions/transfer     & POST & Si  & TransferHandler \\
\end{tablaglab}

El registro (\codigo{server/AuthHandler.java}) recibe \codigo{username} y
\codigo{password}, los valida y delega en \codigo{AuthService}. Devuelve
\codigo{201} si crea la credencial o \codigo{409} si el usuario ya existe.

\begin{lstlisting}[language=Java,caption={Registro de credenciales en AuthHandler}]
private Respuesta registrar(Solicitud s) {
    JsonObject in = parse(s);
    if (in == null) return Respuesta.error(400, "JSON invalido");
    if (!in.has("username") || !in.has("password")) {
        return Respuesta.error(400, "Faltan username/password");
    }
    boolean creado = auth.register(in.get("username").getAsString(),
                                   in.get("password").getAsString());
    if (creado) return Respuesta.json(201, "{\"status\":\"registrado\"}");
    return Respuesta.error(409, "Usuario ya existe");
}
\end{lstlisting}

El login valida las mismas dos llaves y, si las credenciales son correctas,
devuelve el token; si no, responde \codigo{401}.

\begin{lstlisting}[language=none,caption={Respuesta de POST /api/login}]
{ "token": "eyJ0eXAiOiJKV1QiLCJhbGciOiJIUzI1NiJ9.<payload>.<firma>" }
\end{lstlisting}

\subsection{Formatos JSON de cuentas y transferencias}

\codigo{server/AccountHandler.java} extrae el id del ultimo segmento de la ruta
(\codigo{/api/accounts/\{id\}}) sin \codigo{split('/')} (con
\codigo{Long.parseLong} sobre indices, para no alocar un \codigo{String[]} por
peticion), consulta el \codigo{CuentaRepository} y arma el JSON a mano (sin
gson) por ser la ruta caliente. La respuesta lleva \codigo{id},
\codigo{propietario} y \codigo{balance}. Un id no numerico responde
\codigo{400}, pero un id numerico fuera de \codigo{[1, Integer.MAX\_VALUE]} es un
recurso inexistente: responde \codigo{404}.

\begin{lstlisting}[language=Java,caption={Construccion de la respuesta de cuenta}]
long idLargo;
try { idLargo = Long.parseLong(path, ini, fin, 10); }
catch (NumberFormatException e) { return Respuesta.error(400, "Id invalido"); }
if (idLargo < 1 || idLargo > Integer.MAX_VALUE) return Respuesta.error(404, "Cuenta no encontrada");

Cuenta c = repo.get((int) idLargo);
if (c == null) return Respuesta.error(404, "Cuenta no encontrada");

// JSON en un solo StringBuilder: propietario precomputado y saldo desde
// long (sin BigDecimal), para minimizar allocations por GET.
StringBuilder sb = new StringBuilder(96);
sb.append("{\"id\":").append(c.id)
  .append(",\"propietario\":\"").append(c.propietario)
  .append("\",\"balance\":");
c.appendSaldo(sb);
sb.append('}');
return Respuesta.json(200, sb.toString());
\end{lstlisting}

\begin{lstlisting}[language=none,caption={Respuesta de GET /api/accounts/125}]
{ "id": 125, "propietario": "...", "balance": 15750.25 }
\end{lstlisting}

La transferencia (\codigo{server/TransferHandler.java}) espera en el cuerpo
\codigo{sourceAccountId} y \codigo{targetAccountId} como cadenas y
\codigo{amount} como decimal. El monto se convierte a centavos con
\codigo{BigDecimal} para evitar errores de coma flotante antes de invocar
\codigo{CuentaRepository.transferir()}.

\begin{lstlisting}[language=Java,caption={Parseo de ids y monto en TransferHandler}]
origen  = Integer.parseInt(in.get("sourceAccountId").getAsString().trim());
destino = Integer.parseInt(in.get("targetAccountId").getAsString().trim());
centavos = new BigDecimal(in.get("amount").getAsString().trim())
        .movePointRight(2)
        .setScale(0, RoundingMode.HALF_UP)
        .longValueExact();
\end{lstlisting}

\begin{lstlisting}[language=none,caption={Cuerpo de POST /api/transactions/transfer}]
{ "sourceAccountId": "123", "targetAccountId": "456", "amount": 200.00 }
\end{lstlisting}

El resultado del repositorio se traduce a codigos REST: \codigo{200} con
\codigo{seq} si la operacion es correcta, \codigo{404} si una cuenta no existe y
\codigo{400} para misma cuenta, monto invalido o saldo insuficiente.

\begin{tablaglab}{C{1.8cm} L{8.5cm}}
\thead{Codigo} & \thead{Significado en la API} \\ \midrule
200 & Operacion correcta (consulta o transferencia). \\
201 & Credencial registrada. \\
400 & JSON o datos invalidos, faltan campos. \\
401 & Sin JWT valido o credenciales incorrectas. \\
404 & Cuenta no encontrada. \\
409 & Usuario ya existe. \\
\end{tablaglab}

\begin{nota}
\codigo{server/StatsHandler.java} sirve \codigo{GET /stats} sin JWT: es solo
monitoreo de un nodo y entrega el saldo total en centavos para verificar el
invariante de conservacion en una sola peticion.
\end{nota}

\subsection{Autenticacion: JWT y bcrypt}

Las credenciales viven aparte de las cuentas: \codigo{auth/User.java} guarda
\codigo{username} y el hash de la contrasena, y \codigo{auth/UserRepository.java}
las almacena en un \codigo{ConcurrentHashMap} usando \codigo{saveIfAbsent()}
para evitar duplicados. \codigo{auth/AuthService.java} coordina ambos pasos:
hashea al registrar y verifica al iniciar sesion antes de emitir el token.

\begin{lstlisting}[language=Java,caption={Flujo de login en AuthService}]
public String login(String username, String password) {
    User u = repo.findByUsername(username);
    if (u == null) return null;
    if (!PasswordUtil.verify(password, u.getPasswordHash())) return null;
    return JWTUtil.generateToken(username);
}
\end{lstlisting}

El hashing usa bcrypt mediante la libreria jbcrypt en
\codigo{auth/PasswordUtil.java} (\codigo{BCrypt.hashpw} con
\codigo{BCrypt.gensalt} al guardar y \codigo{BCrypt.checkpw} al verificar). Es un
algoritmo deliberadamente lento (~50-100 ms), por eso \codigo{register} y
\codigo{login} corren en el pool WORKER y no en el reactor.

El token se emite y valida con HMAC256 (libreria \codigo{com.auth0} java-jwt) en
\codigo{auth/JWTUtil.java}. El secreto y el issuer son fijos para que un token
emitido por el lider valide en las replicas; el secreto se toma de
\codigo{BANCO\_JWT\_SECRET}.

\begin{lstlisting}[language=Java,caption={Emision y verificacion del JWT}]
public static String generateToken(String username) {
    return JWT.create()
            .withIssuer(ISSUER)
            .withSubject(username)
            .sign(ALGO);
}

public static DecodedJWT verify(String token) throws JWTVerificationException {
    return JWT.require(ALGO).withIssuer(ISSUER).build().verify(token);
}
\end{lstlisting}

\begin{advertencia}
\codigo{BANCO\_JWT\_SECRET} debe ser identico en los tres nodos. Si difiere, un
token valido en el lider sera rechazado con \codigo{401} en las replicas.
\end{advertencia}

\subsection{Validacion del Bearer y cache de tokens}

Los manejadores protegidos llaman a \codigo{HttpUtil.usuarioAutenticado()}
(\codigo{server/HttpUtil.java}) al inicio: si devuelve \codigo{null} responden
\codigo{401}. El metodo lee el header con la clave en minuscula
(\codigo{authorization}, ya normalizada por el parser, para no realocar),
localiza el primer espacio y compara el scheme con \codigo{bearer} de forma
case-insensitive (RFC 7235), recorta el token con \codigo{trim} y verifica la
firma con \codigo{JWTUtil}.

Como la validacion HMAC se repite en cada peticion con el mismo token, los
tokens ya verificados se guardan en un \codigo{ConcurrentHashMap}
(token a subject) con tope de \codigo{TOPE\_CACHE} entradas, evitando recalcular
la firma en la ruta caliente.

\begin{lstlisting}[language=Java,caption={Validacion del Bearer con cache en HttpUtil}]
public static String usuarioAutenticado(Solicitud s) {
    String auth = s.header("authorization");
    if (auth == null) return null;
    int sp = auth.indexOf(' ');
    if (sp < 0) return null;
    if (!auth.substring(0, sp).equalsIgnoreCase("bearer")) return null;
    String token = auth.substring(sp + 1).trim();
    String cacheado = TOKENS_VALIDOS.get(token);
    if (cacheado != null) return cacheado;
    try {
        String sub = JWTUtil.verify(token).getSubject();
        if (TOKENS_VALIDOS.size() < TOPE_CACHE) TOKENS_VALIDOS.put(token, sub);
        return sub;
    } catch (Exception e) {
        return null;
    }
}
\end{lstlisting}

Finalmente, \codigo{server/BancoHttp.java} mapea cada ruta a su manejador y
decide donde se ejecuta. \codigo{/api/accounts} y
\codigo{/api/transactions/transfer} se sirven INLINE en el reactor (son cortas y
de alta frecuencia), mientras que \codigo{/api/register} y \codigo{/api/login}
van al pool WORKER por el costo de bcrypt.

\begin{lstlisting}[language=Java,caption={Mapeo de rutas en BancoHttp}]
Ruteador r = new Ruteador()
        .registrar("/api/register", auth, Ruteador.WORKER)
        .registrar("/api/login", auth, Ruteador.WORKER)
        .registrar("/api/accounts", new AccountHandler(), Ruteador.INLINE)
        .registrar("/api/transactions/transfer", new TransferHandler(), Ruteador.INLINE)
        .registrar("/stats", new StatsHandler(), Ruteador.INLINE)
        // ...
\end{lstlisting}
